\chapter{How an Executor Uses Memory and CPU}
\label{ch:execmem}

\begin{chapterintro}
Spark runs on the Java Virtual Machine: the driver and every executor is a JVM
process, and the engine's memory behaviour, its failure modes, and much of its
internal design follow directly from how the JVM manages memory and runs code.
This chapter builds the memory map of a single executor --- which regions exist,
which a query can use, and which are consumed no matter what --- then adds the
two costs layered on top of it: garbage collection and serialization. By the end
you can size a pod from its configuration, explain a container kill that left no
Java error, and read GC time in the UI. How execution and storage memory borrow
from each other while a job runs is picked up alongside caching in
\Cref{ch:caching}; how to size these regions for a workload is \Cref{ch:tuning}.
\end{chapterintro}

\section{The JVM in one page: bytecode and just-in-time compilation}
\tbd

% To be written alongside whole-stage code generation (\Cref{sec:wholestage-codegen},
% \Cref{ch:tungsten}).

\section{Heap layout and object representation}
\label{sec:mem-heap}

The JVM manages a single contiguous region of memory, the \emph{heap}, from which
every object a program allocates is drawn. A Spark executor is one JVM process,
and the value of \conf{spark.executor.memory} is precisely the maximum size of
that heap. Everything a task does --- decoding a Parquet page into rows, building
a hash table for a join, accumulating a group-by result --- competes for space in
it, and when it fills, the task must either spill to disk or fail.

Two properties of the heap drive most of Spark's decisions about memory. The
first is \emph{per-object overhead}. Every object carries a header --- a mark word
and a pointer to its class --- before any field is stored, and the allocator pads
every object to an eight-byte boundary.

\begin{lstlisting}[numbers=none,caption={Cost of a boxed integer on a 64-bit JVM},captionpos=b]
java.lang.Integer
  object header   12 bytes   (mark word + class pointer)
  int value        4 bytes
  ----------------------
  total           16 bytes   to carry 4 bytes of data
\end{lstlisting}

\noindent A row represented as an array of boxed Java objects (an
\texttt{Integer}, a \texttt{Double}, a \texttt{String} with its own header and
backing character array) can be several times larger in the heap than the same
row encoded as raw bytes. This is the amplification you see when a compressed
Parquet file expands to many times its on-disk size once decoded into Java
objects, and it is why Spark's execution engine stores rows off the object heap
in a packed binary format rather than as JVM objects (\Cref{ch:tungsten}).

The second property is that the heap is \emph{managed}: the program never frees
memory explicitly, and reclamation is the job of the garbage collector
(\Cref{sec:mem-gc}). Spark does not hand the whole heap to a query; it subdivides
it into regions with different owners and different rules, described next.

\section{The executor memory model}
\label{sec:mem-model}

Configuring Spark's memory means deciding how one executor process divides the
memory it is allowed to use. That process holds three separately configured
pools, and confusing them is the most common source of both wasted memory and
unexplained container kills.

\paragraph{The heap.} \conf{spark.executor.memory} sets the maximum size of the
JVM heap, the region every object the executor allocates comes from. Spark carves
it into three parts:

\begin{itemize}
  \item \textbf{Reserved memory} --- a fixed 300~MB Spark holds back for its own
        internal objects. It is not configurable and never available to a query.
  \item \textbf{Unified memory} --- by default 60\% of (heap $-$ 300~MB), set by
        \conf{spark.memory.fraction}. This is the only pool Spark accounts for
        byte by byte and the only one that spills to disk gracefully when it runs
        out. It is shared between \emph{execution memory} (the working memory of a
        running operator: hash tables for joins and aggregations, sort buffers,
        the shuffle write buffer) and \emph{storage memory} (cached DataFrame and
        RDD blocks, and broadcast data held on the executor).
        \conf{spark.memory.storageFraction} (default 0.5) is a soft boundary
        between the two: either side may borrow the other's idle space, but
        execution can always evict borrowed cached blocks to reclaim its half,
        whereas storage cannot push execution below the border. How this plays
        out while a job runs is covered in \Cref{ch:caching}.
  \item \textbf{User memory} --- the remaining 40\%. Everything Spark does not
        manage lives here: the objects your UDFs and RDD closures create, the
        structures a \texttt{mapPartitions} builds up. Spark does not track this
        pool and cannot spill it; overflowing it throws
        \texttt{OutOfMemoryError} and kills the executor.
\end{itemize}

\paragraph{Off-heap memory.} When \conf{spark.memory.offHeap.enabled} is set,
\conf{spark.memory.offHeap.size} adds a second unified-memory pool that lives
outside the heap, addressed directly by the engine and invisible to the garbage
collector (\Cref{sec:mem-gc}). Execution and storage draw on it in addition to
their on-heap pool. It is not taken out of \conf{spark.executor.memory}; it is
additional memory the process reserves, and it must be counted into the
container's limit.

\paragraph{Overhead.} \conf{spark.executor.memoryOverhead} (default the larger
of 384~MB and 10\% of the heap) is memory the process needs that is neither heap
nor off-heap pool: JVM thread stacks, class metadata (Metaspace), buffers for
network transfer and compression codecs, and --- for PySpark --- the Python
worker processes, which can each take hundreds of megabytes. It is not a safety
margin; it is space that is actually consumed, and its default is frequently too
small for PySpark or codec-heavy workloads, in which case the container is killed
with no Spark-level error.

\paragraph{What the cluster manager enforces.} The pod or container is sized as
the sum of all three pools (\Cref{fig:executor-memory-model}). If the process's
real usage crosses that limit --- most often because overhead was underestimated
or user memory grew unbounded --- the kernel kills it
(\Cref{sec:failure-signatures}). The heap merely filling up, by contrast, makes
unified memory spill to disk and only throws when user memory or a single
un-spillable allocation is the cause.

\begin{figure}[htbp]
  \centering
  \includegraphics[width=\linewidth]{executor-memory-model.pdf}
  \caption{An executor's memory: the JVM heap (reserved / unified / user), the
  optional off-heap pool, and the overhead the process needs outside both. The
  cluster manager's container limit must cover all three. Numbers are worked for
  an 8~GB heap with 2~GB overhead.}
  \label{fig:executor-memory-model}
\end{figure}

\paragraph{Per task, not per executor.} Execution memory is shared equally across
the tasks an executor runs at once, which is \conf{spark.executor.cores}. A
larger heap with proportionally more cores gives no single task more room: memory
per task is the execution pool divided by the core count. The number of cores per
executor is therefore a memory decision as much as a parallelism one, and an
executor with too many cores spills on work a narrower one finishes in memory
(\Cref{sec:executor-sizing}).

\section{Garbage collection: Parallel, G1, ZGC}
\label{sec:mem-gc}

Because the heap is managed, the JVM must periodically find objects that are no
longer reachable and reclaim their space. This work is the \emph{garbage
collector} (GC), and for a data-processing workload it is never free: the
collector competes with the application for CPU, and some of its phases are
\emph{stop-the-world}, meaning every application thread in the JVM is frozen until
the phase completes. On an executor a long stop-the-world pause stalls every
running task at once, and if it exceeds the heartbeat timeout the driver declares
the executor lost and re-runs its work elsewhere (\Cref{sec:failure-signatures}).

Collectors differ mainly in how they trade pause length against throughput and
CPU cost.

\begin{itemize}
  \item \textbf{Parallel GC} uses multiple threads but collects the whole heap in
        one stop-the-world pause. It gives the highest raw throughput when pauses
        do not matter, but pause time grows with heap size, so a large executor
        can freeze for seconds at a time.
  \item \textbf{G1} (Garbage-First) divides the heap into equal-sized regions and
        collects a subset of them per cycle, doing most of its marking
        concurrently with the application. Pause times are shorter and roughly
        bounded regardless of heap size, at the cost of more bookkeeping and CPU.
        It is the default at the heap sizes a Spark executor typically runs and
        the collector most Spark tuning assumes.
  \item \textbf{ZGC} does essentially all of its work concurrently and keeps
        pauses under a millisecond even on very large heaps, at a higher constant
        CPU and memory overhead. It is worth choosing only for executors with
        tens of gigabytes of heap where G1 pauses are still a problem.
\end{itemize}

The generational hypothesis --- that most objects die young --- underlies all
three: they collect a small ``young'' area cheaply and often, and the full-heap
collection rarely. A workload that holds a large amount of long-lived data in
cache violates that assumption and pushes the collector toward the expensive
path, which is one more reason to cache deliberately (\Cref{ch:caching}).
Monitoring GC time is covered in \Cref{sec:reading-metrics}; the executor JVM
options that select and tune a collector are in \Cref{ch:tuning}.

\section{Serialization and its cost}
\label{sec:mem-serialization}

The engine keeps a join's hash table or a block of cached rows as one large
off-heap byte buffer rather than millions of small Java objects, which is why
that data is invisible to the garbage collector and costs nothing to collect
(\Cref{sec:mem-model}, \Cref{ch:tungsten}). Getting data into and out of that
flat form is serialization.

\emph{Serialization} is the conversion of an in-memory object graph into a flat
sequence of bytes, and \emph{deserialization} is its inverse. Spark serializes
whenever data leaves a JVM: across the network during a shuffle, to disk when
unified memory spills or a block is cached in serialized form, and when a
broadcast variable is shipped to executors. The cost is both CPU and size.

The two options for arbitrary Java objects differ sharply. The built-in Java
serializer writes the full class name and every field name and type alongside the
values, repeated for every object, so most of the output is self-description
rather than data. \emph{Kryo} writes a compact encoding --- a short numeric class
identifier and the field values --- and produces output several times smaller,
which directly reduces shuffle volume, spill size, and broadcast footprint
(\Cref{ch:tuning}).

DataFrame and SQL workloads sidestep both serializers for most operations: rows
already live in the engine's binary format (\Cref{ch:tungsten}), so moving them
is a memory copy with no object graph to walk. The Java-versus-Kryo choice still
governs RDD code, user-defined functions, and anything that puts arbitrary Java
objects into a shuffle or a broadcast --- the same user-memory objects that
\Cref{sec:mem-model} noted Spark does not manage.

\section{Summary}
\tbd

\exercisehead
\begin{exercises}
\item Given a pod memory limit and a workload description, propose
      \conf{spark.executor.memory}, \conf{spark.executor.cores}, and
      \conf{spark.executor.memoryOverhead}, and justify each number from the
      memory map in \Cref{fig:executor-memory-model}.
\end{exercises}
